\documentclass[11pt]{article}
\usepackage{comment}
\usepackage{graphicx}
\usepackage{fourier-orns}
\usepackage{algorithm, algpseudocode} % For pseudo code
\usepackage{tabto}
\usepackage{amsmath, amsthm, amssymb, amsfonts} % Linear algebra and logic related symbols
\usepackage{color}
\usepackage{fancyhdr} % Header
\usepackage{enumitem} % More option for enumeration and itemization
\usepackage{latexsym}
\usepackage{multicol}
\usepackage{fourier-orns}
\usepackage{algorithm, algpseudocode} % For pseudo code
\usepackage{tabto}
\usepackage{amsmath, amsthm, amssymb, amsfonts}
\usepackage{color}
\usepackage{fancyhdr} % Header
\usepackage{enumitem} % More option for enumeration and itemization
\usepackage{tikz}
\usetikzlibrary{arrows} % graph
\usepackage{url}
\topmargin=-.05in    %0cm
\textheight=9.0in     %24.1cm            
\evensidemargin=0in %0cm
\oddsidemargin=0.0in  %-.3cm
\textwidth=6.5in    %16.5cm 

\usepackage{amsmath}
\usepackage{listings}
\usepackage{xcolor}

\title{MLH PE Fellowship Notes}
\author{Maninder (Kaurman) Kaur}

\begin{document}

\maketitle

2 main databases + use cases:

\begin{itemize}
\item
  Relational databases: data stored in multiple, related tables. Adhere
  to ACID (Atomicity, Consistency, Isolation, and Durability). Usually
  structure data with predefined, fixed relations EXAMPLE: Oracle,
  MySQL, PostgreSQL\\
\item
  NoSQL (Non-Relational): data does not conform to predefined schema.
  Store unstructured data, or semi-structured. EXAMPLE: MongoDB
\end{itemize}

4 types of NoSQL databases:

\begin{itemize}
\item
  Doc-oriented: main unit is a document which is JSON containing
  different fields. Good use case would be good for data analysis and
  looking through records\\
\item
  Columnar: main unit is a column of a table, retrieved column wise.
  Log-base files good use case.\\
\item
  Key-value: all data has a unique key identifier. Good for dealing with
  cache\\
\item
  Graph: main unit contains nodes that represent entities and edges that
  represent relationships. Good from graph like structures and social
  networks.
\end{itemize}

Database management system (DBMS)\\
SQL (Structured Query Language) used on Relational DBMS

Cardinality - distinctiveness of info values contained in a column. High
cardinality means column contains an outsized proportion of distinctive
values. Low cardinality means more repeats.\\
Mapping constraint is a data constraint that express the \# of entities
to which another entity can be related via a relationship set. For
binary relationship set R on an entity A and B, there are four possible
mapping cardinalities:

\begin{itemize}
\item
  One to one 1:1\\
\item
  One to many 1:M\\
\item
  Many to one M:1\\
\item
  Many to Many M:M
\end{itemize}

Service - group of processes (known as daemons) running continuously in
background.Most Linux distros uses same process manager: systemd

\texttt{Systemd} is a system + service manager for Linux, allows
parallel startup of system services at boot time, on-demand activation
of daemons, and dependency based on service control logic.
\texttt{.service} -\textgreater{} system services, \texttt{.target}
-\textgreater{} group of systemd units. Main features:

\begin{itemize}
\item
  Socket-based activation allows parallel services to start without
  losing messages when restarted.\\
\item
  Starts services in parallel as soon as all listening sockets are in
  place. Significantly reduces time required to boot system\\
\item
  Only stops running services on shutdown
\end{itemize}

\texttt{systemctl} - command used to control and manage \texttt{systemd}
services\\
Activate - \texttt{systemctl\ start\ foo}, deactivate -
\texttt{systemctl\ stop\ foo}, restart -
\texttt{systemctl\ restart\ foo}, view status -
\texttt{systemctl\ status\ foo}. More commands -\textgreater{}
\texttt{systemctl\ status\ foo}

HTTP (hypertext transfer protocol) - set of rules for sending and
receiving web pages\\
Web request - communicative message that goes between client/web browser
to servers.\\
HTTP Communication happens between client and server. Client makes HTTP
request and gets back a HTTP response.

HTTP request line consists of Method, URL and Version\\
METHODS:

\begin{itemize}
\item
  GET: used to seek info from specified source.\\
\item
  HEAD: only returns headers\\
\item
  POST: seeks to create and make a new object\\
\item
  PUT: replace fully with new section

  URL: location that HTTP is referring to\\
  VERSION: what version? HTTP/3? HTTP/2? HTTP/1.1?\\
  HTTP request is followed by HTTP Header Structure:
\item
  General: applied on request and response w/o affecting database body\\
\item
  Request: contains info about fetched\\
\item
  Response: contains location of source requested\\
\item
  Entity: info about the body of resources
\end{itemize}

HTTP response has initial status line, some header lines and an entity
body\\
Status Codes:

\begin{itemize}
\item
  1xx: informational\\
\item
  2xx: success (200 is request was successful)\\
\item
  3xx: redirection\\
\item
  4xx: client error (404 is file not found, 400 is bad request/not in
  format system can read)\\
\item
  5xx: server error (500 is server error, 505 is requested HTTP version
  not supported)
\end{itemize}

Command: \texttt{curl\ http://example.org\ –head\ -silent}: curl is name
of command, --head or -I tells curl to send an HTTP request with head
method. And -silent tells curl to not display progress.

Caching - supplements primary database by removing unnecessary pressure
on it. Different types of caching:

\begin{itemize}
\item
  Database Integrated: managed within database engine and has built in
  write through capabilities. The cons are size and capabilities.\\
\item
  Local: stored within application, removing network traffic. Con is
  that each app has its own node, thus it cannot shared with other local
  caches.\\
\item
  Remote: stored in dedicated servers by NoSQL.
\end{itemize}

\end{document}
